% Overleaf entry point: compile with XeLaTeX; latexmk runs Biber automatically.
% Compile this file with XeLaTeX and Biber; references.bib is the bibliography.
\documentclass[10pt,a4paper]{article}
\usepackage[left=1in,right=1in,top=1.25in,bottom=1.25in]{geometry}
\usepackage{fontspec}
\IfFontExistsTF{Times New Roman}{\setmainfont{Times New Roman}}{\setmainfont{TeX Gyre Termes}}
\IfFontExistsTF{Arial}{\setsansfont{Arial}}{\setsansfont{TeX Gyre Heros}}
\usepackage{amsmath,booktabs,tabularx,array,graphicx}
\usepackage{titlesec}
\usepackage[american]{babel}
\usepackage{csquotes}
\usepackage[backend=biber,style=apa,natbib=true,maxbibnames=20,minbibnames=19,maxcitenames=2,mincitenames=1,eprint=false]{biblatex}
\DeclareLanguageMapping{american}{american-apa}
\addbibresource{references.bib}
\setlength{\bibitemsep}{0pt}
\setlength{\bibparsep}{0pt}
\setlength{\bibhang}{0.5in}
\renewcommand*{\bibfont}{\fontsize{10}{12}\selectfont}
\usepackage[hidelinks]{hyperref}
\urlstyle{same}
\setlength{\parindent}{0.5in}
\setlength{\parskip}{0pt}
\setlength{\textfloatsep}{9pt}
\setlength{\intextsep}{9pt}
\titleformat{\section}{\sffamily\bfseries\fontsize{12}{14}\selectfont}{}{0pt}{}
\titleformat{\subsection}{\sffamily\fontsize{12}{14}\selectfont}{}{0pt}{}
\titlespacing*{\section}{0pt}{9pt}{3pt}
\titlespacing*{\subsection}{0pt}{6pt}{3pt}
\pagestyle{empty}
\setlength{\emergencystretch}{2em}
\newcommand{\supporturl}{https://drive.google.com/drive/folders/1BjS3toqoV5py_ZhA1HjCcs162TQ6f7-t}
\begin{document}
\begin{center}
{\sffamily\bfseries\fontsize{14}{16}\selectfont What Mathematical Work Do AI Tutors Leave for Learners?\\A Qualitative Study and a Test of Decision-Guided Tutoring\par}
\vspace{8pt}
Jingying He, Jiawei Ni, Peilin Ying, Duc Le, and Shayan Ahmad\\
University of Pennsylvania
\end{center}
\vspace{2pt}
\begin{quotation}\noindent
\textbf{Abstract:} AI tutors must help students without completing the mathematical reasoning they need to practice. Using TutorMoments-Preview, a dataset of real tutoring conversations, we examine this problem in two ways. First, we code the work left for learners in 38 moments. Explicit instructions about balancing support and challenge produced more replies inviting connections, explanations, or strategies than general tutoring guidance. Second, we design an experiment in which a decision model chooses scaffolding or rigor and a language model writes the tutor's reply. Across 520 moments, we compare two decision systems, Jev and OpenAI Decisions, with the tutor choosing its own approach and with fixed instructions. Automated first-reply pass rates were 68.7\% for self-choice, 72.9\% with Jev, 73.5\% with Decisions, and 74.0\% with rigor instructions at every moment. Decision-guided replies had higher pass rates than self-choice, but neither system outperformed fixed rigor guidance.

\end{quotation}


\section{Introduction}
A mathematics tutor must decide when to help and when to leave the next step to the learner. \emph{Scaffolding} makes a task accessible, for example by offering a hint or a useful representation. \emph{Pushing for rigor} asks the learner to do more demanding mathematical work, such as explaining a relationship or choosing a strategy. Too much help can remove the reasoning that the student needs to practice. Too little can leave the student stuck. This assistance dilemma is central to the design of tutoring systems \citep{koedinger2007assistance}. Research on teacher questioning also shows why the work elicited from students matters, beyond whether the teacher's explanation is correct \citep{franke2009questioning}.

We study this problem using TutorMoments-Preview, a public collection of real tutoring conversations and AI-generated continuations \citep{zhang2026tutormoments}. Our project has two connected parts. First, we examine what mathematical work remains for a learner after a tutor responds. Second, we test a technical approach intended to improve that response: a separate model reads the conversation, chooses scaffolding or rigor, and selects instructions for the model that writes the tutor's reply. The idea is that a small tutor might respond more appropriately if it receives a clear teaching decision before generating text.

The approach builds on work that separates instructional planning from language generation. StratL guides tutors through teaching intents \citep{puech2025stratl}, while DICE separates decisions about whether to intervene from the ability to produce an answer \citep{pal2026dice}. Eduardo takes a different approach by training smaller tutors through reinforcement learning \citep{macina2026eduardo}. We investigate whether an additional decision step can help an existing model without retraining it. Our questions are: \textbf{(1)} How do tutoring instructions change the mathematical work left for learners? \textbf{(2)} Does a separate scaffold-or-rigor decision improve a smaller model's reply beyond detailed instructions alone, and what does that decision step add in cost and delay?

\section{What TutorMoments contains and how we use it}
TutorMoments-Preview contains 462 de-identified transcripts of remote tutoring with students in grades 2--7. A transcript records dialogue and descriptions of shared-screen activity. A \emph{moment} is a stopping point in a conversation where a tutor must decide what to do next. Experienced teachers described important situations, the human tutor's action, and what followed. The dataset authors used language models to map these descriptions to categories and combined annotations by majority vote. Their balanced evaluation set contains 520 moments: 260 labeled as opportunities for scaffolding and 260 as opportunities to push for rigor \citep{zhang2026tutormoments}.

For each moment, the benchmark gives an AI tutor the conversation up to that point and asks it to continue with a simulated learner. It compares two kinds of instructions. The \emph{plain prompt} gives general tutoring guidance. The \emph{evaluation-aware prompt} explains scaffolding, rigor, and over-scaffolding, and asks the tutor to choose appropriately. It does not reveal the moment's assigned category. The authors found that explicit guidance improved performance, although models still tended to give too much help and miss opportunities for greater challenge. Their scoring can credit a response that combines scaffolding and rigor, provided it avoids over-scaffolding. The categories therefore identify instructional opportunities, not two mutually exclusive forms of good teaching.

Our qualitative study uses the released human and AI continuations. Our technical study generates new tutor replies from the same conversation histories. Neither study collects new classroom data. We examine the first response after each stopping point, rather than reproducing the benchmark's multi-turn exchange. Our percentages consequently measure a different outcome from its published leaderboard.

\section{Study 1: What work does the tutor leave for the learner?}
We compared the human continuation with GPT-5.5 continuations under the plain and evaluation-aware prompts at 38 mathematics moments, one per student and 19 per source category. The sample began with 40 moments selected from 226 eligible situations after excluding students used in earlier pilots. Two were removed because intervening learner activity made the human response an unsuitable comparison. They were not replaced. A separate audit established a 346-moment mathematics comparison set with clear conversation boundaries and usable responses across the original conditions. The supplement records the complete selection process and exclusions.

The unit of analysis is the first tutor response block, combining consecutive tutor utterances until the learner next speaks or acts. Following a theory-informed content-analysis approach \citep{hsieh2005content}, we developed two codes. \emph{Work Left} describes the most demanding mathematical contribution still available to the learner: no new work or repetition (W0), a specified calculation or heavily cued recognition (W1), a connection or explanation within a narrowed path (W2), or a strategy, argument, or diagnosis with meaningful discretion (W3). \emph{Information Supplied} describes how much of the current answer the tutor provides, from none (R0), through general guidance (R1) and a task-specific step or relationship (R2), to the complete answer (R3). Unclear cases receive NA. A higher W code does not necessarily mean better teaching: the appropriate demand depends on the learner.

Two AI-assisted coding passes applied the codebook to responses with their source identities hidden and preceding dialogue available. We report the first pass and retain the second to show how interpretation changes the results. All 114 response texts and their supporting quotations were checked against the source packets. The main comparison is between 38 pairs of AI replies. Human continuations provide a naturalistic reference because they were not produced under controlled prompt conditions.

\subsection{Qualitative findings}
Evaluation-aware replies more often left a connection, explanation, or strategy for the learner: 16/38 received W2 or W3, compared with 2/38 plain replies (Table~\ref{tab:work}). Plain replies supplied the complete active answer in 21/38 cases, compared with 5/38 evaluation-aware replies. One plain reply was unclear on this measure.

\begin{table}[ht]
\centering\small
\caption{Mathematical work left in the first coding pass ($n=38$ per column).}
\label{tab:work}
\begin{tabularx}{\linewidth}{@{}Xrrr@{}}
\toprule
Work left & Plain & Evaluation-aware & Human reference\\
\midrule
W0: No new work or repetition &21&2&20\\
W1: Calculation or cued recognition &15&20&12\\
W2: Connection or explanation &2&14&1\\
W3: Strategy, argument, or diagnosis &0&2&0\\
NA: Unclear demand &0&0&5\\
\bottomrule
\end{tabularx}
\end{table}

For example, one problem asked the learner to complete $25\times14=25\times(2\times\star)$. The plain reply stated, ``So the star should be 7,'' then verified the equality. The evaluation-aware reply asked, ``What number times 2 makes 14?'' The latter still supplied the decisive relationship, but left the calculation to the student. This distinction is easy to miss if analysis only records whether the tutor asked a question. The second coding pass also favored evaluation-aware replies, assigning W2/W3 to 23 rather than five plain replies. Its different counts show that boundaries between calculation and mathematical explanation require judgment.

\section{Study 2: Can a separate decision improve the next tutor reply?}
\subsection{The proposed approach}
We attempted a technical contribution by separating two tasks normally handled by one language model: deciding what kind of support is needed and writing that support. A \emph{decision model}, also called a router here, reads the conversation and returns either scaffolding or rigor. That choice selects one of two detailed instruction sets for the tutor. Scaffolding instructions ask for the smallest useful hint while leaving the key calculation or inference to the student. Rigor instructions ask the tutor to invite a justification, connection, or independent step without supplying the reasoning first. In both cases, the tutor still receives the full conversation and must write a reply appropriate to the current problem.

We tested two decision systems: Jev 1.14.0, a TypeSafe model designed to make fast structured choices \citep{almeida2026jev}, and OpenAI's Decisions API using GPT-6 Luna \citep{openai2026decisions}. Qwen3-32B, an open-weight language model, wrote the tutor replies \citep{qwen2025model}. The intended sequence is therefore: \textit{conversation so far $\rightarrow$ scaffold-or-rigor decision $\rightarrow$ corresponding instructions $\rightarrow$ Qwen tutor reply}. The motivation is practical: specialized decision making could improve a less expensive tutor, but it is only useful if the resulting replies justify the extra step.

\subsection{Separating decision accuracy from response quality}
We first tested whether each decision system matched the dataset's assigned category across all 520 moments. The selected prompts produced 376 matches for Jev (72.3\%) and 386 for Decisions (74.2\%). These predictions were saved and reused in the response experiment. Prompt development had already used this dataset, so these are exploratory results on previously inspected cases. A subsequent prompt revision reduced agreement for both systems. The supplement preserves all versions.

We then tested what those choices meant for the actual tutor reply. At every moment, Qwen generated three alternatives: \textbf{self-choice}, where it received both detailed instruction sets and chose an approach itself; \textbf{scaffold-guided}, where it received the scaffolding instructions; and \textbf{rigor-guided}, where it received the rigor instructions. All three also received the publisher's evaluation-aware prompt. This produced 1,560 replies. For each saved Jev or Decisions prediction, we selected the corresponding scaffold-guided or rigor-guided reply. Generating both alternatives let us compare the decision systems on exactly the same candidate text, without confounding their choices with different generations.

We included two simple controls: always use scaffold guidance and always use rigor guidance. These test whether a separate decision model adds value beyond a consistent instruction. A further diagnostic selects the reply corresponding to the dataset label, showing what happens when that category is supplied correctly. This label is never provided to either decision model. All models receive the complete preceding dialogue, with future turns and teacher annotations excluded. The tutor additionally receives the publisher's student metadata. Generation used temperature zero, non-thinking mode, and a 384-token output limit. No reply was truncated.

\subsection{How replies were evaluated}
Gemini 2.5 Flash evaluated the three replies in randomized order, without condition names or dataset labels. A reply passed when it displayed the behavior called for by the dataset label, avoided over-helping and mathematical errors, and fit the conversation. A reply judged to combine scaffolding and rigor could satisfy either category. We also asked Claude Sonnet 4.6 to evaluate the same replies for a fixed subset of 32 students, balanced by category. Comparisons pair replies from the same moments, with uncertainty intervals obtained by resampling students so that related moments remain together.

\subsection{Technical findings}
Detailed self-choice passed 357/520 cases (68.7\%), Jev-directed selection 379 (72.9\%), and Decisions-directed selection 382 (73.5\%; Table~\ref{tab:routing}). Jev's paired gain was 4.2 percentage points (95\% student-cluster bootstrap interval: $-1.0$ to $9.2$); Decisions' gain was 4.8 points ($0.2$ to $9.5$). Neither exceeded fixed rigor guidance (385/520, 74.0\%). Jev improved rigor-condition passes from 149 to 180 but reduced scaffolding-condition passes from 208 to 199. Thus, its aggregate improvement concealed a trade-off across situations.

\begin{table}[ht]
\centering\small
\caption{Passing first replies under the primary evaluator.}
\label{tab:routing}
\begin{tabularx}{\linewidth}{@{}Xrrrr@{}}
\toprule
Policy & All 520 & Scaffold 260 & Rigor 260 & Matched 346\\
\midrule
Detailed self-choice &68.7\%&80.0\%&57.3\%&72.0\%\\
Jev-directed &72.9\%&76.5\%&69.2\%&72.8\%\\
Decisions-directed &73.5\%&80.8\%&66.2\%&75.1\%\\
Always scaffold &56.5\%&86.2\%&26.9\%&59.8\%\\
Always rigor &74.0\%&65.8\%&82.3\%&76.0\%\\
Reference-label diagnostic &84.2\%&86.2\%&82.3\%&86.4\%\\
Opposite-label diagnostic &46.3\%&65.8\%&26.9\%&49.4\%\\
\bottomrule
\end{tabularx}
\vspace{3pt}
\begin{minipage}{\linewidth}\footnotesize
The last column uses the strict matched mathematics frame. The two label-based policies are diagnostic controls, not deployable systems. All policies share the same paired scaffold/rigor replies, except self-choice.
\end{minipage}
\end{table}

The stricter mathematics comparison was less favorable to routing: Jev passed 252/346 cases versus 249 for self-choice, a 0.9-point difference ($-5.3$ to $6.5$). Decisions passed 260 and fixed rigor 263. On the 32-case evaluation audit, Claude passed 16 self-choice replies and 22 each for Jev, Decisions, and fixed rigor. On those same cases, the primary evaluator passed 20 self-choice and 19 Jev replies, reversing the relative result. Overall agreement was 75/96 replies (78.1\%): condition differences were evaluator-sensitive.

Category agreement and reply success diverged. When Jev matched the reference category, 320/376 selected replies passed; when it disagreed, 59/144 still passed. Across the paired policy replies, both passed at 214 moments, only rigor passed at 171, only scaffolding at 80, and neither at 55. Even reference-label selection passed only 438/520. Better category prediction therefore leaves a separate response-execution problem, while category disagreement does not invariably produce a poor reply. In moment 24, for example, the learner was working on $1{,}000-10$. Both routers correctly selected scaffolding, but the selected reply returned to an earlier $338+586$ problem; the evaluator marked a context failure.

Median generator time to first answer was 1.20 seconds for self-choice. Adding archived routing time yielded reconstructed medians of 1.43 seconds for Jev and 1.60 for Decisions. Estimated selected-response costs per 1,000 cases were \$0.59, \$0.90, and \$1.25, respectively. These estimates exclude evaluation and do not establish a deployment speedup. The complete new experiment cost approximately \$3.00, including evaluation and conservative failed-call reservations.


\section{Discussion and contribution}
The qualitative analysis describes what a learner can contribute after a reply. The technical analysis tests whether a separate decision helps produce that reply. Together they show why choosing an instructional category is only part of the tutoring problem. A model can choose scaffolding correctly yet supply too much information or address the wrong task. Conversely, a response can combine support and challenge even when its assigned instruction names only one of them.

Our technical results provide limited support for adding a binary router. Both learned selectors improved the full-preview point estimate relative to self-choice, but fixed rigor guidance matched or exceeded them, and Jev's gain nearly disappeared in the strict matched mathematics frame. This does not imply that tutors should always demand rigor. The fixed-rigor replies also expressed scaffolding in many cases, and the evaluator recognized both behaviors in 202/520 replies under that policy. A policy name is therefore an imperfect description of the support actually delivered. The useful design target may be selecting a specific next demand and the information to withhold, rather than treating scaffolding and rigor as mutually exclusive response types. The paired branch analysis makes this distinction inspectable without attributing differences to separate sampled replies.


These findings come from selected benchmark moments, previously inspected data, AI-assisted coding, and automated judgments. They do not measure student learning. The human replies are observations from real sessions, while the new AI replies were not delivered to learners. Independent human evaluation and testing on new conversations are needed before drawing conclusions about classroom benefit. The contribution of this course study is a concrete coding account of mathematical work and a reproducible, controlled test showing where decision-guided tutoring helps, where simple guidance performs similarly, and why category agreement alone is insufficient evidence of a better tutor.

\section{Supplementary materials}
The \href{https://drive.google.com/drive/folders/1PRhd04y0frYdvlDwIcqd27XcJX7hLSpD}{qualitative analysis folder} provides the codebook, coded responses, results, and source conversations. The \href{https://drive.google.com/drive/folders/1Oqy731wJQkNdRJABpMrL6EE4DM9BPZ5B}{technical analysis folder} provides the full methods, prompts, predictions, replies, ratings, and runnable code. A \href{https://github.com/shayanahmad7/tutormoments-decision-router}{GitHub repository} maintains the implementation; the Drive copy is accessible without repository access. AI tools assisted coding, programming, evaluation, and drafting. No independent human coding validation is claimed.

\clearpage
\printbibliography[heading=bibliography,title={References}]
\end{document}
